% --- Archon physics formalization source begin ---
% archon:physics
% archon:covers PhyXMiniProblems/problem_phyx_mini_0473.lean
% archon:source-report reports/phyx_mini/problem_phyx_mini_0473.source.json

\chapter{Physics problem phyx\_mini\_0473}
\label{ch:PhyXMiniProblems_problem_phyx_mini_0473}

\paragraph{Problem source.}
\#\# Physical scenario

A gasoline truck engine takes in 10,000 J of heat and delivers 2000 J of mechanical work per cycle. The heat is obtained by burning gasoline with heat of combustion L\_c = 5.0 \textbackslash{}times 10\textasciicircum{}4 J/g.

\#\# Question

How much gasoline is burned in each cycle?

\#\# Answer choices

- A: A: 0.25g

- B: B: 0.20g

- C: C: 0.30g

- D: D: 0.35g

\#\# Figure caption (auxiliary; use the image as primary evidence)

The image is a diagram representing a thermodynamic process. It includes:

1. **Vertical Arrows**: 
   - An arrow pointing downward labeled with \textbackslash{}( T\_H \textbackslash{}) at the top and \textbackslash{}( T\_C \textbackslash{}) at the bottom, indicating a temperature gradient from high (\textbackslash{}( T\_H \textbackslash{})) to low (\textbackslash{}( T\_C \textbackslash{})).
   
2. **Horizontal Arrow**: 
   - An arrow pointing to the right labeled \textbackslash{}( W = 2000 \textbackslash{}, \textbackslash{}text\{J\} \textbackslash{}), indicating the work done by the system.

3. **Text and Labels**:
   - \textbackslash{}( Q\_H = 10,000 \textbackslash{}, \textbackslash{}text\{J\} \textbackslash{}) next to the downward arrow, representing the heat input at the higher temperature.
   - \textbackslash{}( Q\_C = ? \textbackslash{}) next to the downward arrow near \textbackslash{}( T\_C \textbackslash{}), indicating the heat expelled at the lower temperature, which is unknown.

The diagram depicts an energy balance concept, often used in the context of heat engines or refrigerators.

\#\# Dataset metadata

Laws of Thermodynamics; Physical Model Grounding Reasoning, Multi-Formula Reasoning

\paragraph{Recorded answer/context.}
B: B: 0.20g

\paragraph{Figure/image path.}
/root/proposal\_for\_physic/science-mango/phyx\_mini\_run/phyx\_data/test\_image/473.png

\paragraph{Formalization target.}
create a compiling Lean file with sorry bodies at `PhyXMiniProblems/problem\_phyx\_mini\_0473.lean`. The Lean declarations must preserve the physical quantities, dimensions or dimensional roles, figure labels, governing-law hypotheses, and final relation expressed by this problem.
Use Mathlib/Physlib names found through LeanExplore where available. If a physics API is missing, introduce faithful local abstractions rather than scalar placeholder aliases.

\begin{theorem}[Physics formalization target]
\label{thm:physics:phyx_mini_0473:target}
\lean{PhyXMiniProblems.ProblemPhyXMini0473.gasoline_burned_per_cycle_eq_point_two_grams}
\uses{def:physics:phyx-mini-0473:phyxminiproblems-problemphyxmini0473-fuelmassingrams, def:physics:phyx-mini-0473:phyxminiproblems-problemphyxmini0473-gasolinetruckenginecycle, def:physics:phyx-mini-0473:phyxminiproblems-problemphyxmini0473-matchesproblemstatementandprimaryfigure, def:physics:phyx-mini-0473:phyxminiproblems-problemphyxmini0473-obeysgasolineheatenginecyclelaws}
One engine cycle burns $0.20$ grams of gasoline.
\end{theorem}
\begin{proof}
The figure and problem readouts give heat input $Q_H=10000$ joules and combustion energy $H=5\cdot10^4$ joules per gram.  The combustion law says $Q_H=mH$.  Substitute the two readouts and divide by the positive value $5\cdot10^4$ to obtain $m=10000/(5\cdot10^4)=1/5=0.20$ grams.
\end{proof}

% --- Archon declaration topology begin ---
\section{Declaration topology}
\begin{definition}[Fuel Mass Quantity]
  \label{def:physics:phyx-mini-0473:phyxminiproblems-problemphyxmini0473-fuelmassquantity}
  \lean{PhyXMiniProblems.ProblemPhyXMini0473.FuelMassQuantity}
  Nonnegative gasoline mass with physical dimension M
\end{definition}
\begin{proof}
  This is the declared model definition of Fuel Mass Quantity.
\end{proof}
\begin{definition}[Specific Combustion Energy Quantity]
  \label{def:physics:phyx-mini-0473:phyxminiproblems-problemphyxmini0473-specificcombustionenergyquantity}
  \lean{PhyXMiniProblems.ProblemPhyXMini0473.SpecificCombustionEnergyQuantity}
  Heat released per unit fuel mass. Its dimension is energy divided by mass, namely L² T⁻²; its SI readout is in joules per kilogram
\end{definition}
\begin{proof}
  This is the declared model definition of Specific Combustion Energy Quantity.
\end{proof}
\begin{definition}[energy In Joules]
  \label{def:physics:phyx-mini-0473:phyxminiproblems-problemphyxmini0473-energyinjoules}
  \lean{PhyXMiniProblems.ProblemPhyXMini0473.energyInJoules}
  Read a dimensionful energy magnitude in joules
\end{definition}
\begin{proof}
  This is the declared model definition of energy In Joules.
\end{proof}
\begin{definition}[fuel Mass In Grams]
  \label{def:physics:phyx-mini-0473:phyxminiproblems-problemphyxmini0473-fuelmassingrams}
  \lean{PhyXMiniProblems.ProblemPhyXMini0473.fuelMassInGrams}
  \uses{def:physics:phyx-mini-0473:phyxminiproblems-problemphyxmini0473-fuelmassquantity}
  Read a physical fuel mass in grams
\end{definition}
\begin{proof}
  This is the declared model definition of fuel Mass In Grams.
\end{proof}
\begin{definition}[combustion Energy In Joules Per Gram]
  \label{def:physics:phyx-mini-0473:phyxminiproblems-problemphyxmini0473-combustionenergyinjoulespergram}
  \lean{PhyXMiniProblems.ProblemPhyXMini0473.combustionEnergyInJoulesPerGram}
  \uses{def:physics:phyx-mini-0473:phyxminiproblems-problemphyxmini0473-specificcombustionenergyquantity}
  Read a specific combustion energy in joules per gram. A coherent-SI readout of this dimension is in joules per kilogram, and 1 J/g = 1000 J/kg
\end{definition}
\begin{proof}
  This is the declared model definition of combustion Energy In Joules Per Gram.
\end{proof}
\begin{definition}[Figure Arrow Direction]
  \label{def:physics:phyx-mini-0473:phyxminiproblems-problemphyxmini0473-figurearrowdirection}
  \lean{PhyXMiniProblems.ProblemPhyXMini0473.FigureArrowDirection}
  Arrow orientations appearing in the heat-engine diagram
\end{definition}
\begin{proof}
  This is the declared model definition of Figure Arrow Direction.
\end{proof}
\begin{definition}[Figure Temperature Label]
  \label{def:physics:phyx-mini-0473:phyxminiproblems-problemphyxmini0473-figuretemperaturelabel}
  \lean{PhyXMiniProblems.ProblemPhyXMini0473.FigureTemperatureLabel}
  The two literal reservoir-temperature labels in the diagram
\end{definition}
\begin{proof}
  This is the declared model definition of Figure Temperature Label.
\end{proof}
\begin{definition}[Figure Energy Label]
  \label{def:physics:phyx-mini-0473:phyxminiproblems-problemphyxmini0473-figureenergylabel}
  \lean{PhyXMiniProblems.ProblemPhyXMini0473.FigureEnergyLabel}
  The three literal energy-flow labels in the diagram
\end{definition}
\begin{proof}
  This is the declared model definition of Figure Energy Label.
\end{proof}
\begin{definition}[Supplied Heat Engine Figure]
  \label{def:physics:phyx-mini-0473:phyxminiproblems-problemphyxmini0473-suppliedheatenginefigure}
  \lean{PhyXMiniProblems.ProblemPhyXMini0473.SuppliedHeatEngineFigure}
  \uses{def:physics:phyx-mini-0473:phyxminiproblems-problemphyxmini0473-figurearrowdirection, def:physics:phyx-mini-0473:phyxminiproblems-problemphyxmini0473-figuretemperaturelabel, def:physics:phyx-mini-0473:phyxminiproblems-problemphyxmini0473-figureenergylabel}
  Data transcribed from primary image 473.png. The two numeric fields are display readouts, while the cold-side heat is explicitly marked unknown
\end{definition}
\begin{proof}
  This is the declared model definition of Supplied Heat Engine Figure.
\end{proof}
\begin{definition}[Gasoline Truck Engine Cycle]
  \label{def:physics:phyx-mini-0473:phyxminiproblems-problemphyxmini0473-gasolinetruckenginecycle}
  \lean{PhyXMiniProblems.ProblemPhyXMini0473.GasolineTruckEngineCycle}
  \uses{def:physics:phyx-mini-0473:phyxminiproblems-problemphyxmini0473-fuelmassquantity, def:physics:phyx-mini-0473:phyxminiproblems-problemphyxmini0473-specificcombustionenergyquantity, def:physics:phyx-mini-0473:phyxminiproblems-problemphyxmini0473-suppliedheatenginefigure}
  Independent physical observables for one truck-engine cycle. In particular, the gasoline mass is not defined from an answer choice, and the rejected heat is retained even though the question asks only for fuel consumption
\end{definition}
\begin{proof}
  This is the declared model definition of Gasoline Truck Engine Cycle.
\end{proof}
\begin{definition}[Matches Problem Statement And Primary Figure]
  \label{def:physics:phyx-mini-0473:phyxminiproblems-problemphyxmini0473-matchesproblemstatementandprimaryfigure}
  \lean{PhyXMiniProblems.ProblemPhyXMini0473.MatchesProblemStatementAndPrimaryFigure}
  \uses{def:physics:phyx-mini-0473:phyxminiproblems-problemphyxmini0473-energyinjoules, def:physics:phyx-mini-0473:phyxminiproblems-problemphyxmini0473-combustionenergyinjoulespergram, def:physics:phyx-mini-0473:phyxminiproblems-problemphyxmini0473-gasolinetruckenginecycle}
  The numerical problem statement and qualitative primary-image transcription. No field specifies the gasoline mass or any answer-choice value
\end{definition}
\begin{proof}
  This is the declared model definition of Matches Problem Statement And Primary Figure.
\end{proof}
\begin{definition}[Obeys Gasoline Heat Engine Cycle Laws]
  \label{def:physics:phyx-mini-0473:phyxminiproblems-problemphyxmini0473-obeysgasolineheatenginecyclelaws}
  \lean{PhyXMiniProblems.ProblemPhyXMini0473.ObeysGasolineHeatEngineCycleLaws}
  \uses{def:physics:phyx-mini-0473:phyxminiproblems-problemphyxmini0473-energyinjoules, def:physics:phyx-mini-0473:phyxminiproblems-problemphyxmini0473-fuelmassingrams, def:physics:phyx-mini-0473:phyxminiproblems-problemphyxmini0473-combustionenergyinjoulespergram, def:physics:phyx-mini-0473:phyxminiproblems-problemphyxmini0473-gasolinetruckenginecycle}
  The first law for a cyclic heat engine and the fuel-combustion energy law. These are general relations among independent observables, not numerical assertions of the requested fuel mass
\end{definition}
\begin{proof}
  This is the declared model definition of Obeys Gasoline Heat Engine Cycle Laws.
\end{proof}
% --- Archon declaration topology end ---
% --- Archon physics formalization source end ---
